\slajdid{10-s058}
\begin{frame}[label=10-s058]{Totem pol}
\gusttekst\setlength{\abovedisplayskip}{3pt}\setlength{\belowdisplayskip}{3pt}\setlength{\abovedisplayshortskip}{2pt}\setlength{\belowdisplayshortskip}{2pt}\begin{columns}[c,onlytextwidth]\column{.42\textwidth}\centering\input{slike/tikz/s058_totem_pol_neopterecen.tex}\column{.1\textwidth}\centering\tikz\draw[->,DEplava,>=Latex](0,0)--(1.3,0);\column{.44\textwidth}\centering\begin{tikzpicture}[x=.85cm,y=.55cm,font=\sitneoznake,>=Latex]\ctikzset{bipoles/length=.6cm}\draw(0,-.5)--(0,.5);\draw(-1,0)--(0,0);\draw(-1,0)to[R,l=$R_B$](-1,2.5);\draw(0,.32)--(.7,.7)--(.7,1)to[R,l_=$R_C$](.7,2.5);\draw(-1.3,2.5)--(1,2.5);\node[above]at(-.5,2.5){$V_{CC}$};\draw[->](0,-.32)--(.7,-.7);\node[right]at(.2,0){T1};\draw(.7,-.7)--(.7,-1.1);\draw(.45,-1.1)--(.95,-1.1)--(.7,-1.65)--cycle;\draw(.45,-1.65)--(.95,-1.65);\draw(.7,-1.65)--(.7,-2)--(1.6,-2)node[ocirc]{}node[right]{$V_o$};\node[right]at(1,-1.35){D1};\draw(1.1,-2)to[R,l=$R\to\infty$](1.1,-3.6)node[ground]{};\end{tikzpicture}
\end{columns}
\par\medskip Kako otpornost u emitoru teži beskonačnosti
\[I_E\to0\ \Rightarrow\ I_C\to0\land I_C\to0\ \Rightarrow\ V_C\to V_{CC}\land V_B\to V_{CC}\]
Napon između baze i kolektora je jednak nuli, tranzistor je u aktivnom režimu. Radi sa jako malim strujama pa možemo reći da je na ivici provođenja. Isto važi i za diodu D1. U tom slučaju
\[V_0=V_{CC}-R_B I_{B1}-V_{BE1}-V_{D1}=V_{CC}-0-V_{\gamma T}-V_{\gamma D}=V_{CC}-V_{\gamma T}-V_{\gamma D}\]

\end{frame}
